\chapter{BNP (B-Type Natriuretic Peptide)}

\section*{What Is This Test?}
B-type natriuretic peptide (BNP, also called brain natriuretic peptide) is a 32-amino-acid cyclic peptide hormone produced primarily by ventricular cardiomyocytes in response to increased ventricular wall stress from volume expansion and pressure overload. BNP is synthesized as a 134-amino-acid precursor (pre-proBNP), which is cleaved to proBNP (108 amino acids) and then further cleaved by the serine proteases corin and furin to yield the biologically active BNP (32 amino acids) and an inactive N-terminal fragment, NT-proBNP (76 amino acids). BNP binds to natriuretic peptide receptor A (NPR-A) in the kidneys, vasculature, adrenal glands, and heart, increasing intracellular cGMP and producing: natriuresis (renal sodium excretion), diuresis, vasodilation (both arterial and venous), inhibition of the renin-angiotensin-aldosterone system (decreased renin and aldosterone), inhibition of sympathetic nervous system activity, and antifibrotic and antihypertrophic effects on the myocardium. These counter-regulatory effects oppose the volume expansion and vasoconstriction of heart failure. BNP and NT-proBNP are the two clinically used natriuretic peptide assays, both reflecting the same pathophysiological process but with different characteristics: BNP (half-life 20 minutes) is cleared by the natriuretic peptide C receptor (NPR-C) and neutral endopeptidase (neprilysin); NT-proBNP (half-life 60--120 minutes) is cleared primarily by the kidneys and accumulates in renal failure. Both are used for diagnosis and exclusion of acute heart failure in the emergency setting, risk stratification, and monitoring chronic heart failure. BNP and NT-proBNP are the most important cardiac biomarkers alongside troponin.

\section*{Alternative Names}
B-Type Natriuretic Peptide; BNP; Brain Natriuretic Peptide (historical name, never produced in brain); NT-proBNP (N-Terminal pro-B-Type Natriuretic Peptide); Serum BNP; Plasma BNP; Natriuretic Peptide Type B; ProBNP; NT-proBNP by Immunoassay

\section*{Principle}
BNP and NT-proBNP are measured using two-site sandwich immunometric (chemiluminescent) immunoassays. For BNP: two monoclonal antibodies against different epitopes on the BNP-32 molecule (typically the ring structure and the C-terminal region) are used. For NT-proBNP: antibodies against the N-terminal portion of proBNP are used. Common platforms: Abbott Architect (BNP), Roche Cobas ECLIA (NT-proBNP), Siemens Atellica (BNP and NT-proBNP), Beckman Coulter (BNP). The assays are NOT standardized to each other, and BNP and NT-proBNP values are NOT interchangeable (conversion is not possible because of different half-lives and clearance mechanisms) \cite{tietz2012textbook}. Specimens should be collected into EDTA (purple-top) tubes for BNP (BNP is rapidly degraded by serine proteases; EDTA inhibits this degradation). NT-proBNP is more stable and can be collected into either EDTA or serum separator tubes. Results are available in 15--30 minutes on point-of-care platforms (i-STAT, Triage) and within 1 hour in central laboratories.

\section*{History}
BNP was first isolated from porcine brain extracts in 1988 by Sudoh and colleagues, leading to the misnomer "brain" natriuretic peptide; it was subsequently found to be produced primarily by the heart. The first radioimmunoassay for BNP was developed in the early 1990s. The landmark Breathing Not Properly study (Maisel et al., 2002, N Engl J Med) demonstrated that BNP measurement at the point of care accurately distinguished acute heart failure from other causes of dyspnea in the emergency department (BNP >100 pg/mL had 90\% sensitivity and 76\% specificity; BNP >500 pg/mL had 90--95\% positive predictive value) \cite{maisel2002rapid}. The PRIDE study (Januzzi et al., 2005) established NT-proBNP cutoffs for heart failure diagnosis (age-stratified: <50 years >450 pg/mL; 50--75 years >900 pg/mL; >75 years >1800 pg/mL) \cite{januzzi2005nterminal}. The development of sacubitril/valsartan (Entresto, approved 2015), a neprilysin inhibitor that increases BNP levels by inhibiting its degradation, introduced new complexity to BNP interpretation in treated patients. The 2017 ACC/AHA/HFSA focused update and 2021 ESC heart failure guidelines cemented BNP/NT-proBNP as essential biomarkers in heart failure diagnosis and management \cite{mcdonagh2021esc}.

\section*{Why This Test Is Ordered}
\begin{itemize}
  \item Diagnosis of acute decompensated heart failure (ADHF) in patients presenting to the emergency department with acute dyspnea --- BNP >100 pg/mL (or NT-proBNP >300 pg/mL) supports the diagnosis; BNP <100 pg/mL (NT-proBNP <300 pg/mL) has a 90--98\% negative predictive value and effectively excludes ADHF \cite{ponikowski2016esc}
  \item Excluding heart failure as the cause of dyspnea in the outpatient setting --- BNP <35 pg/mL or NT-proBNP <125 pg/mL effectively rules out chronic heart failure
  \item Risk stratification in acute coronary syndromes --- elevated BNP/NT-proBNP predicts increased risk of death, heart failure, and recurrent MI independent of troponin and other risk factors
  \item Prognostic assessment in chronic heart failure --- rising BNP/NT-proBNP predicts hospitalization and death; BNP >200 pg/mL or NT-proBNP >1000 pg/mL identifies high-risk patients
  \item Monitoring response to heart failure therapy --- titration of guideline-directed medical therapy (ACE inhibitors/ARBs/ARNI, beta-blockers, mineralocorticoid receptor antagonists, SGLT2 inhibitors) guided by BNP/NT-proBNP trends
  \item Screening for left ventricular dysfunction in high-risk populations (post-MI, diabetes, hypertension, known coronary artery disease)
  \item Assessment of volume status in patients with chronic kidney disease and suspected volume overload
  \item Differentiating cardiac from pulmonary causes of dyspnea when clinical assessment is equivocal
  \item Pre-operative cardiac risk stratification in non-cardiac surgery
  \item Evaluating treatment response to sacubitril/valsartan (Entresto) --- BNP increases during therapy (neprilysin inhibition decreases BNP degradation) while NT-proBNP decreases (reflects reduced cardiac wall stress) --- NT-proBNP is the preferred marker during Entresto therapy
\end{itemize}

\section*{Primary vs Secondary Test}
BNP/NT-proBNP is a primary diagnostic test for acute heart failure in the emergency department and a primary prognostic marker in heart failure. The ACC/AHA guidelines recommend BNP/NT-proBNP measurement for diagnosis of heart failure when the diagnosis is uncertain, for establishing prognosis, and for guiding therapy \cite{yancy2017acc}. It is the most widely used biomarker in cardiology alongside high-sensitivity cardiac troponin.

\section*{How This Test Is Performed}
A healthcare professional draws a blood sample from a vein into an EDTA (purple-top) tube for BNP (sodium or lithium heparin tubes are acceptable on some platforms; check with your laboratory). For NT-proBNP, serum separator tube (gold-top) is also acceptable. For BNP, the specimen should be centrifuged within 4 hours and plasma separated because BNP is degraded by circulating proteases at room temperature. NT-proBNP is more stable (24 hours at room temperature, 7 days refrigerated). The patient should be seated or supine for 10--15 minutes before venipuncture, as upright posture increases BNP/NT-proBNP by 10--15\%. The test can be performed at the point of care (whole blood, fingerstick, or venous) with results in 15--20 minutes, or in the central laboratory with results in 1 hour.

\section*{What Preparation Is Needed}
No special preparation is typically required in the emergency setting. In elective settings, fasting is not required. The patient should rest for 10--15 minutes before venipuncture. Strenuous exercise should be avoided for 2--3 hours before testing because exercise transiently elevates BNP/NT-proBNP (myocardial stretch). High-dose biotin (>5 mg/day) must be held 48 hours before testing. In patients on sacubitril/valsartan (Entresto, a neprilysin inhibitor), BNP is pharmacologically elevated because neprilysin degrades BNP; NT-proBNP (not degraded by neprilysin) is the preferred biomarker during Entresto therapy. Note the timing of blood draw relative to the last Entresto dose (trough BNP or NT-proBNP is most informative).

\section*{Contraindications and Precautions}
Factors that elevate BNP/NT-proBNP independent of heart failure: (1) Age: BNP and NT-proBNP increase progressively with age (age-stratified cutoffs are essential for NT-proBNP); (2) Female sex: baseline BNP/NT-proBNP is 10--20\% higher in women than men (estrogen effects); (3) Chronic kidney disease (CKD): both BNP and NT-proBNP are elevated because of decreased renal clearance; NT-proBNP is more dependent on renal clearance and is more affected by CKD than BNP; (4) Atrial fibrillation: atrial stretch and myocyte stress elevate BNP/NT-proBNP independently of ventricular dysfunction; (5) Acute coronary syndromes: myocardial ischemia causes transient BNP elevation even without heart failure; (6) Right ventricular dysfunction and pulmonary hypertension (including pulmonary embolism); (7) Sepsis, systemic inflammatory response syndrome (SIRS), and critical illness; (8) Hyperthyroidism; (9) Cirrhosis with ascites; (10) Anemia; (11) Subarachnoid hemorrhage and stroke; (12) High cardiac output states (pregnancy, hyperthyroidism). Factors that lower BNP/NT-proBNP: (1) Obesity --- BMI inversely correlated with BNP/NT-proBNP, possibly due to increased NPR-C clearance receptors on adipocytes; obese patients with heart failure may have "falsely low" BNP levels. Lower BNP cutoffs should be used in obesity (BNP <50 pg/mL to exclude heart failure if BMI >35). (2) Pericardial constriction and restrictive cardiomyopathy may have relatively low BNP compared to degree of heart failure because the pericardium or infiltrated myocardium limits wall stretch, which is the stimulus for BNP production. (3) Flash pulmonary edema: BNP may not have time to rise because symptom onset is rapid (<1--2 hours). (4) Sacubitril/valsartan therapy: increases BNP (neprilysin inhibition decreases BNP degradation) but decreases NT-proBNP (reduced wall stress from improved heart function). Always interpret BNP with knowledge of the patient's medications. NT-proBNP is preferred in patients on Entresto. Hemolysis can interfere with BNP assays.

\section*{Risks}
Minimal risk. Slight pain, bruising, or hematoma at venipuncture site. Rarely: vasovagal syncope. No radiation exposure.

\section*{What Happens After the Test}
Results are available within 15--30 minutes (point of care) or 1 hour (central lab). BNP <100 pg/mL effectively excludes acute heart failure (NPV 90--98\%) and alternative causes of dyspnea should be sought (COPD exacerbation, asthma, pneumonia, pulmonary embolism, anxiety). BNP 100--500 pg/mL is consistent with heart failure but also seen in right ventricular dysfunction, pulmonary embolism, and renal failure; clinical correlation and echocardiography are needed. BNP >500 pg/mL strongly suggests acute decompensated heart failure (PPV >90\%). NT-proBNP age-stratified cutoffs: <50 years: >450 pg/mL; 50--75 years: >900 pg/mL; >75 years: >1800 pg/mL supports acute heart failure. NT-proBNP <300 pg/mL effectively excludes acute heart failure. In the outpatient setting: BNP <35 pg/mL or NT-proBNP <125 pg/mL excludes chronic heart failure. In patients with confirmed HF, BNP/NT-proBNP levels should decrease by 30--50\% with effective therapy; failure to decrease is an adverse prognostic sign and should prompt more aggressive treatment \cite{mueller2019heart}.

\section*{Understanding Results}

\subsection*{Below Normal}
\begin{itemize}
  \item BNP <100 pg/mL (or NT-proBNP <300 pg/mL) in acute dyspnea: Effectively excludes acute decompensated heart failure as the cause of dyspnea (NPV 90--98\%). Alternative diagnoses should be pursued: COPD exacerbation, asthma exacerbation, pneumonia, pulmonary embolism, pneumothorax, pleural effusion, or anxiety/hyperventilation.
  \item BNP <35 pg/mL (or NT-proBNP <125 pg/mL) in the outpatient setting: Excludes chronic heart failure (NPV 95--98\%). Symptoms are unlikely due to heart failure. Screening for other causes is appropriate.
  \item Decreasing BNP/NT-proBNP during heart failure treatment: A decrease of >30--50\% from admission to discharge indicates effective decongestion with diuretics and indicates favorable prognosis. Failure to decrease predicts increased re-hospitalization and mortality.
  \item Obesity (BMI >35): BNP/NT-proBNP may be "falsely low" relative to the degree of heart failure. Use lower cutoffs: BNP <50 pg/mL to exclude HF in severely obese patients.
  \item Flash pulmonary edema: BNP may be <100--200 pg/mL despite florid pulmonary edema because the time course (minutes) is too short for significant BNP synthesis and release.
\end{itemize}

\subsection*{Above Normal}
\begin{itemize}
  \item Acute decompensated heart failure (ADHF): BNP >500--1000 pg/mL or NT-proBNP >1800--4500 pg/mL (age-dependent). The degree of elevation correlates with left ventricular filling pressures (wedge pressure) and the severity of heart failure. BNP >1000 pg/mL and NT-proBNP >5000 pg/mL identify very high-risk patients with 30-day mortality exceeding 10--15\%. Clinical context: dyspnea, orthopnea, paroxysmal nocturnal dyspnea, elevated JVP, crackles, S3 gallop, peripheral edema.
  \item Chronic heart failure (stable): BNP 100--500 pg/mL or NT-proBNP 300--1800 pg/mL. Levels reflect ongoing ventricular wall stress and correlate with disease severity and prognosis. Serial monitoring every 3--6 months is recommended.
  \item Acute coronary syndrome (ACS): myocardial ischemia transiently elevates BNP/NT-proBNP even in the absence of clinical heart failure. Elevated BNP/NT-proBNP post-MI is a powerful independent predictor of mortality and subsequent heart failure, superior to LVEF in some studies.
  \item Atrial fibrillation: atrial myocyte stretch elevates BNP/NT-proBNP. Levels increase by 20--50\% during AF compared to sinus rhythm. Rate control or cardioversion reduces BNP levels.
  \item Chronic kidney disease (CKD): BNP and NT-proBNP are both elevated in CKD due to decreased renal clearance. NT-proBNP is more affected (clearance is >90\% renal) than BNP (clearance via NPR-C and neprilysin, ~50\% renal). In dialysis patients, NT-proBNP is markedly elevated and BNP is moderately elevated even without heart failure. Higher diagnostic cutoffs are needed (BNP >200--300 pg/mL for HF in CKD).
  \item Pulmonary embolism and pulmonary hypertension: Right ventricular strain increases BNP/NT-proBNP. Elevated BNP in acute PE identifies patients at increased risk of adverse outcomes and may guide thrombolysis decisions.
  \item Sepsis and critical illness: systemic inflammation, cytokine release, and volume resuscitation (causing ventricular stretch) elevate BNP/NT-proBNP. Elevated BNP in sepsis predicts myocardial dysfunction and mortality.
  \item Entresto (sacubitril/valsartan) therapy: BNP increases by 20--50\% after initiation because neprilysin inhibition prevents BNP degradation, while NT-proBNP decreases by 20--40\% because cardiac function improves and wall stress is reduced. NT-proBNP is the preferred biomarker for monitoring Entresto-treated patients because it reflects the true reduction in cardiac wall stress.
  \item Advanced age: BNP/NT-proBNP increase with age. Age-stratified cutoffs are essential, particularly for NT-proBNP.
\end{itemize}

\section*{Normal Results}
\begin{itemize}
  \item BNP (acute dyspnea, emergency setting): <100 pg/mL (excludes acute heart failure; NPV >90\%)
  \item BNP (outpatient, chronic): <35 pg/mL (excludes chronic heart failure)
  \item BNP gray zone: 100--500 pg/mL (consistent with heart failure but also seen in RV dysfunction, PE, CKD, AF)
  \item BNP likely ADHF: >500 pg/mL (PPV >90\%)
  \item NT-proBNP (acute dyspnea): <300 pg/mL (excludes acute HF; age-independent cut-off)
  \item NT-proBNP (age-stratified for HF diagnosis): Age <50: >450 pg/mL; Age 50--75: >900 pg/mL; Age >75: >1800 pg/mL (supports acute HF)
  \item NT-proBNP (outpatient, chronic): <125 pg/mL (excludes chronic heart failure)
  \item NT-proBNP likely ADHF: >1000--1800 pg/mL (age-dependent)
  \item BNP and NT-proBNP reference ranges are assay-specific and NOT interchangeable. Always use the cutoffs for the specific assay being performed.
  \item Note: BNP and NT-proBNP values are affected by numerous non-cardiac factors. Always interpret in the clinical context of symptoms, physical examination findings, chest X-ray, and echocardiography.
\end{itemize}

\section*{Follow-up Tests}
\begin{itemize}
  \item Echocardiography (transthoracic, 2D with Doppler): Essential for confirming the diagnosis of heart failure and determining the ejection fraction (HFrEF, HFmrEF, HFpEF). Assesses left ventricular size, wall thickness, systolic/diastolic function, valvular disease, and right ventricular function. The cornerstone of heart failure evaluation.
  \item Chest X-ray: Evaluates for pulmonary edema (cephalization, Kerley B lines, pleural effusions), cardiomegaly, and alternative causes of dyspnea (pneumonia, pneumothorax).
  \item High-sensitivity cardiac troponin (hs-cTn): Evaluate for concurrent myocardial injury or acute coronary syndrome causing elevated BNP. Prognostic in heart failure.
  \item Comprehensive metabolic panel (electrolytes, BUN, creatinine, GFR): Renal function is essential for interpreting BNP/NT-proBNP. Hyponatremia is a marker of severe heart failure.
  \item Complete blood count: Exclude anemia as a contributor to high-output heart failure and dyspnea.
  \item TSH: Exclude hyperthyroidism (elevates BNP, can cause high-output heart failure) or hypothyroidism contributing to dyspnea.
  \item 12-lead electrocardiogram (ECG): Evaluate for ischemia, arrhythmia (atrial fibrillation), left ventricular hypertrophy, and QRS duration (for CRT candidacy).
  \item Liver function tests: Hepatic congestion in right heart failure.
  \item Serum albumin: Volume status and nutritional assessment. Low albumin contributes to pulmonary edema.
  \item Right heart catheterization (Swan-Ganz): Gold standard for measuring filling pressures and cardiac output. Indicated when diagnosis is uncertain (BNP elevation without clear echocardiographic abnormality) or for cardiogenic shock.
  \item Serial BNP/NT-proBNP: Every 1--3 months during heart failure management to assess treatment response. A decrease of >30--50\% indicates effective therapy.
  \item Cardiorenal syndrome monitoring: Serial BNP and creatinine to balance decongestion with renal function preservation.
\end{itemize}

\section*{References}
\bibliographystyle{unsrtnat}
\bibliography{references}

\clearpage
